\subsubsection{多尺度粗糙读出：连续极限与不可导可见量}\label{subsubsec:groupoid-zeckendorf-multiscale-rough-readout}

注 \ref{rem:residual} 已将残差解释为值纤维到规范横截面之间仍需外置承载的自由度；命题 \ref{prop:fold-renyi-collision-identity} 及附录 \ref{subsubsec:fold-fiber-groundstate-freezing-oracle-capacity} 则把 \(d_m(x)\)、\(S_q(m)\) 与预言机容量 \(\mathcal C_m(B)\) 固定为可审计的纤维尺度。本段讨论与之互补的一条读出链：若这些残差自由度不完全停留在独立账本中，而是沿分辨率层级以确定性方式重新注入可见读出，则可见层已经足以出现连续但不可导的几何。由此，不可微性不再需要被解释为本体随机，而可被解释为多尺度确定性读出的结果。

\begin{definition}[残差相容的多尺度可见读出]\label{def:groupoid-zeckendorf-multiscale-rough-readout}
令
\[
\Omega_\infty:=\{0,1\}^{\NN},
\qquad
\tau_{\infty\to m}:\Omega_\infty\to \Omega_m
\]
为取前 \(m\) 位前缀的截断映射。对每个 \(m\ge 1\) 与 \(x\in X_m\)，固定一个纤维编号
\[
\iota_{m,x}:\Fold_m^{-1}(x)\xrightarrow{\sim}\{1,\dots,d_m(x)\}.
\]
记
\[
D_m:=\max_{x\in X_m}d_m(x),
\qquad
r_m(\omega):=
\iota_{m,\Fold_m(\tau_{\infty\to m}(\omega))}
\bigl(\tau_{\infty\to m}(\omega)\bigr)
\in \{1,\dots,d_m(\Fold_m(\tau_{\infty\to m}(\omega)))\}.
\]

再取一列振幅 \(a_m>0\)、频率 \(\lambda_m>0\)、连续轮廓 \(\psi_m:\RR/2\ZZ\to\RR\) 及相位映射
\[
\Theta_m:\{1,\dots,D_m\}\to \RR/2\ZZ.
\]
若
\[
\sum_{m\ge 1} a_m\|\psi_m\|_\infty<\infty,
\]
则定义第 \(M\) 阶截断读出
\[
Y_\omega^{(\le M)}(t):=
\sum_{m=1}^{M}
a_m\,
\psi_m\!\bigl(\lambda_m t+\Theta_m(r_m(\omega))\bigr),
\qquad t\in\RR,
\]
并把一致极限
\[
Y_\omega(t):=\lim_{M\to\infty}Y_\omega^{(\le M)}(t)
\]
称为由 Fold--fiber 残差诱导的多尺度可见量；上述按 \(r_m(\omega)\) 因子化的层级读出即称为残差相容。
\end{definition}

\begin{proposition}[连续极限的存在]\label{prop:groupoid-zeckendorf-multiscale-rough-readout-continuity}
在定义 \ref{def:groupoid-zeckendorf-multiscale-rough-readout} 的假设下，对每个 \(\omega\in\Omega_\infty\)，级数
\[
\sum_{m\ge 1}
a_m\,
\psi_m\!\bigl(\lambda_m t+\Theta_m(r_m(\omega))\bigr)
\]
在 \(\RR\) 上一致收敛，从而 \(Y_\omega\) 为连续函数。
\leanverified{paper\_groupoid\_zeckendorf\_multiscale\_rough\_readout\_continuity}
\end{proposition}
\begin{proof}
对任意 \(t\in\RR\) 与 \(m\ge 1\)，有
\[
\left|
a_m\,
\psi_m\!\bigl(\lambda_m t+\Theta_m(r_m(\omega))\bigr)
\right|
\le
a_m\|\psi_m\|_\infty.
\]
由假设 \(\sum_{m\ge 1} a_m\|\psi_m\|_\infty<\infty\)，Weierstrass 判别法给出一致收敛。连续函数级数的一致极限仍连续，故 \(Y_\omega\) 连续。证毕。
\end{proof}

\begin{definition}[局部振荡、点态 H\"older 指数与 \(p\)-variation]\label{def:groupoid-zeckendorf-visible-roughness-certificates}
设 \(Y:I\to\RR\) 为连续函数，\(t\in I\) 为内点。
\begin{enumerate}
  \item 对 \(\delta>0\)，定义局部振荡
  \[
  \operatorname{osc}(Y;t,\delta):=
  \sup_{|s-t|\le \delta}Y(s)-\inf_{|s-t|\le \delta}Y(s).
  \]
  \item 定义 \(Y\) 在 \(t\) 的点态 H\"older 指数为
  \[
  h_Y(t):=
  \sup\Bigl\{
  \alpha>0:\ \exists C,\delta_0>0,\ 
  |Y(s)-Y(t)|\le C|s-t|^\alpha
  \ \forall |s-t|\le \delta_0
  \Bigr\}.
  \]
  \item 对闭区间 \(J=[u,v]\subset I\) 与 \(p>0\)，定义 \(p\)-variation
  \[
  V_p(Y;J):=
  \sup_{\Pi}
  \sum_{j=1}^{n}
  |Y(t_j)-Y(t_{j-1})|^p,
  \]
  其中上确界取遍 \(J\) 的一切有限分割
  \(
  \Pi:u=t_0<t_1<\cdots<t_n=v
  \)。
\end{enumerate}
\end{definition}

\begin{proposition}[局部 H\"older 阈值与可导性的必要条件]\label{prop:groupoid-zeckendorf-visible-holder-nondifferentiability}
若连续函数 \(Y\) 在内点 \(t\) 可导，则
\[
\operatorname{osc}(Y;t,\delta)=O(\delta)
\qquad (\delta\downarrow 0),
\]
并且 \(h_Y(t)\ge 1\)。特别地，若存在 \(\alpha<1\)、常数 \(c>0\) 与 \(\delta_n\downarrow 0\) 使
\[
\operatorname{osc}(Y;t,\delta_n)\ge c\,\delta_n^\alpha,
\]
则 \(Y\) 在 \(t\) 不可导。
\end{proposition}
\begin{proof}
若 \(Y\) 在 \(t\) 可导，则
\[
Y(s)-Y(t)=Y'(t)(s-t)+o(|s-t|)
\qquad (s\to t).
\]
因此存在常数 \(C,\delta_0>0\) 使
\[
|Y(s)-Y(t)|\le C|s-t|
\qquad (|s-t|\le \delta_0).
\]
于是
\[
\operatorname{osc}(Y;t,\delta)
\le
2\sup_{|s-t|\le \delta}|Y(s)-Y(t)|
\le 2C\delta
\qquad (0<\delta\le \delta_0),
\]
并由点态 H\"older 指数的定义得到 \(h_Y(t)\ge 1\)。

反过来，若存在 \(\alpha<1\) 与序列 \(\delta_n\downarrow 0\) 使
\(
\operatorname{osc}(Y;t,\delta_n)\ge c\,\delta_n^\alpha
\),
则
\[
\frac{\operatorname{osc}(Y;t,\delta_n)}{\delta_n}
\ge
c\,\delta_n^{\alpha-1}\xrightarrow[n\to\infty]{}\infty.
\]
这与可导时 \(\operatorname{osc}(Y;t,\delta)=O(\delta)\) 矛盾，故 \(Y\) 在 \(t\) 不可导。证毕。
\end{proof}
\leanverified{paper\_groupoid\_zeckendorf\_visible\_holder\_nondifferentiability}

\begin{proposition}[平滑投影的不可导边界]\label{prop:groupoid-zeckendorf-smooth-projection-nogo}
设 \(M\) 为有限维 \(C^1\) 流形，\(\gamma:I\to M\) 为 \(C^1\) 轨道，\(O:M\to\RR\) 为 \(C^1\) 读出。则
\[
Y(t):=O(\gamma(t))
\]
仍为 \(C^1\) 函数。特别地，有限维平滑隐藏轨道经平滑读出不会自动生成处处不可导的可见量。
\end{proposition}
\leanverified{paper\_groupoid\_zeckendorf\_smooth\_projection\_nogo}
\begin{proof}
由链式法则，
\[
Y'(t)=DO_{\gamma(t)}\bigl(\gamma'(t)\bigr),
\]
右端关于 \(t\) 连续，故 \(Y\in C^1(I)\)。证毕。
\end{proof}

\begin{remark}[粗糙可见量的来源边界]\label{rem:groupoid-zeckendorf-smooth-projection-boundary}
命题 \ref{prop:groupoid-zeckendorf-smooth-projection-nogo} 表明，“高维到低维的普通投影”不足以解释处处不可导的可见曲线。若粗糙性确实出现，则其来源至少应包含下列机制之一：无限层级的相位再注入、多尺度频率放大、粗粒化极限，或非有限维状态空间中的确定性动力学。Fold--fiber 口径自然对应前两类机制，因为残差标签本身就是沿分辨率层级组织的。
\end{remark}

\begin{theorem}[确定性多尺度读出的不可导实现]\label{thm:groupoid-zeckendorf-deterministic-rough-visible-quantity}
取奇整数 \(b\ge 3\) 与常数 \(a\in(0,1)\)，满足
\[
ab>1+\frac{3\pi}{2}.
\]
\leanverified{paper\_groupoid\_zeckendorf\_deterministic\_rough\_visible\_quantity}
令
\[
M:=(\RR/2\ZZ)^{\NN},
\qquad
\Phi_t\bigl((\theta_m)_{m\ge 1}\bigr):=
(\theta_m+b^m t)_{m\ge 1},
\]
并定义读出
\[
O\bigl((\theta_m)_{m\ge 1}\bigr):=
\sum_{m\ge 1} a^m\cos(\pi\theta_m).
\]
则 \(O\) 在 \(M\) 上定义良好且连续；对基点 \(0=(0,0,\dots)\in M\)，对应的可见量
\[
Y(t):=O(\Phi_t(0))
=
\sum_{m\ge 1} a^m\cos(\pi b^m t)
\]
连续且处处不可导。特别地，确定性动力学配合多尺度读出已足以产生连续但不可导的可见层。
\end{theorem}
\begin{proof}
对任意 \((\theta_m)_{m\ge 1}\in M\)，有
\[
\left|a^m\cos(\pi\theta_m)\right|\le a^m,
\qquad
\sum_{m\ge 1} a^m<\infty.
\]
故 \(O\) 由连续函数级数一致收敛定义，因而连续。于是 \(Y(t)=O(\Phi_t(0))\) 连续。

另一方面，令 \(\psi_m=\cos(\pi\cdot)\)、\(\lambda_m=b^m\)、\(\Theta_m\equiv 0\)，则 \(Y\) 正是定义 \ref{def:groupoid-zeckendorf-multiscale-rough-readout} 中残差相容读出的特例。其处处不可导性由经典 Weierstrass--Hardy 判据给出：当 \(b\) 为奇整数且 \(ab>1+\frac{3\pi}{2}\) 时，级数
\[
\sum_{m\ge 1} a^m\cos(\pi b^m t)
\]
连续且处处不可导。证毕。
\end{proof}

\begin{proposition}[网格差商的指数发散证书]\label{prop:groupoid-zeckendorf-rough-visible-difference-quotient-certificate}
\leanverified{paper\_groupoid\_zeckendorf\_rough\_visible\_difference\_quotient\_certificate}
在定理 \ref{thm:groupoid-zeckendorf-deterministic-rough-visible-quantity} 的设定下，令
\[
h_M:=b^{-M}\qquad (M\ge 1).
\]
则
\[
\left|
\frac{Y(h_M)-Y(0)}{h_M}
\right|
\ge
\frac{2a}{1-a}(ab)^{M-1}.
\]
因此沿网格 \(h_M=b^{-M}\) 的差商绝对值至少以 \((ab)^M\) 的指数尺度发散；这给出一个有限分辨率可复核的粗糙证书。
\end{proposition}
\begin{proof}
由定义，
\[
\frac{Y(h_M)-Y(0)}{h_M}
=
b^M\sum_{m\ge 1} a^m\bigl(\cos(\pi b^{m-M})-1\bigr).
\]
当 \(m\ge M\) 时，\(b^{m-M}\) 为奇整数，故
\[
\cos(\pi b^{m-M})-1=-2.
\]
当 \(m<M\) 时，\(\cos(\pi b^{m-M})-1\le 0\)。于是
\[
\left|
\frac{Y(h_M)-Y(0)}{h_M}
\right|
\ge
2b^M\sum_{m\ge M}a^m
=
2b^M\frac{a^M}{1-a}
=
\frac{2a}{1-a}(ab)^{M-1}.
\]
证毕。
\end{proof}

\begin{remark}[残差相位与碰撞矩接口]\label{rem:groupoid-zeckendorf-rough-readout-fiber-interface}
定理 \ref{thm:groupoid-zeckendorf-deterministic-rough-visible-quantity} 只说明粗糙可见量在确定性制度下确实存在；它所用的 \(\Theta_m\equiv 0\) 仍是相位冻结的特例。真正与 Fold--fiber 主链相接的情形，是让 \(\Theta_m\) 对残差标签 \(r_m(\omega)\) 产生非平凡分离。此时，纤维多重度 \(d_m(x)\)、碰撞矩 \(S_q(m)\)、冻结相阈值（定理 \ref{thm:fold-pressure-freezing-threshold}）、预言机容量（定理 \ref{thm:fold-oracle-capacity-closed-form}）以及最小在线延迟 \(\delta_{\min}(\Fold_\infty)=3\)（定理 \ref{thm:online-delay-fold} 与命题 \ref{prop:pom-delay-min}）共同给出可见粗糙度的几何、统计与时间约束。换言之，纤维残差并不只控制可逆恢复所需的侧信息预算，也控制可见层能够被抹平到何种程度。
\end{remark}
